\documentclass[fleqn]{article} 
% hmargin is set narrow (0.7cm) so that each of the 2 columns has
% the same width (~272pt) as the original A3-landscape 4-column
% layout. This keeps the wide 6x6-matrix equation from overflowing
% and reproduces the original line-wrapping / page count (2 pages).
% includeheadfoot reserves room for the running header/footer.
\usepackage[a4paper,vmargin=0.8cm,hmargin=0.7cm,includeheadfoot,headsep=4pt,footskip=14pt]{geometry}

\usepackage{latexsym}
\usepackage{graphicx}
\usepackage{latexsym}

\usepackage{multicol}
\usepackage{fancyhdr}
\usepackage{lastpage}

\usepackage{epsfig}
\usepackage[dvips,xdvi]{color}

\usepackage{amsmath}
\usepackage{amssymb}

%--------------------------------------------------------------
% textwidth/textheight/margins are controlled by the geometry
% package above (a4paper). The hardcoded A3-landscape values
% below were removed so the layout fits A4 portrait.
%	\setlength{\textwidth}{1120pt}
%	\setlength{\textheight}{790pt}
%--------------------------------------------------------------
%	\setlength{\oddsidemargin}{-50pt}
%	\setlength{\evensidemargin}{-50pt}
%--------------------------------------------------------------
%	\setlength{\topmargin}{-70pt}
	\setlength{\headheight}{15pt}
	\setlength{\topskip}{0pt}
%--------------------------------------------------------------
	\setlength{\columnseprule}{0.5pt}
	\setlength{\columnsep}{10pt}
%--------------------------------------------------------------
	\renewcommand{\baselinestretch}{1}
%--------------------------------------------------------------
	\setlength{\mathindent}{0.5cm}
	\setlength{\parindent}{7pt}
%--------------------------------------------------------------
	\renewcommand{\bottomfraction}{0.99}
	\renewcommand{\topfraction}{0.99}
	\renewcommand{\textfraction}{0.01}
%--------------------------------------------------------------
%	\pagestyle{empty}
%--------------------------------------------------------------


%----------------------------------------------------------------
%	definition
%----------------------------------------------------------------
%----------------------------------------------------------------
%   equation 環境
%----------------------------------------------------------------
%		\def\topmat{\ \\ \vspace{-4mm}}
%		\def\botmat{\ \\ \vspace{-2mm}}
%		\def\bes{\vspace*{2mm}}
%		\def\ees{\ \\[-2mm]}
	%........................................................
	% equationの環境設定(\be,\ee)
	%      例 \bea 式 & 式 & 式 \nonumber\\ 式 & 式 & 式 \eea
	%........................................................
%		\def\be{\\ \ \\[-12mm] \begin{equation}}
%		\def\ee{\end{equation} \ \\[-11mm]}
		\newcommand{\be}{\begin{equation}}
		\newcommand{\ee}{\end{equation}}
	%........................................................
	% eqnarrayの環境設定(\bea,\eea)
	%      例 \bea 式 & 式 & 式 \nonumber\\ 式 & 式 & 式 \eea
	%........................................................
%		\def\bea{\\ \ \\[-16mm] \begin{eqnarray}}
%		\def\eea{\end{eqnarray} \ \\[-13mm]}
		\newcommand{\bea}{\begin{eqnarray}}
		\newcommand{\eea}{\end{eqnarray}}
	%........................................................
	% 番号なしeqnの環境設定(\nbe,\nee)
	%      例 \nbe 式 \\ 式  \nee
	%........................................................
%		\def\nbe{\\ \ \\[-12mm] \[}
%		\def\nee{\] \ \\[-11mm]}
		\newcommand{\nbe}{\[}
		\newcommand{\nee}{\]}
	%........................................................
	% matrixの入力(\bmatrix,\ematrix)
	%	例 \bmatrix{cc} 1&2\\3&4 \ematrix
	%........................................................
		\renewcommand{\bmatrix}{\left[\begin{array}}
		\newcommand{\ematrix}{\end{array}\right]}
	%........................................................
	% Greek character
	%........................................................
		\def\gggm{\mu}
		\def\ggge{\varepsilon}
		\def\gggw{\omega}
		\def\gggh{\theta}
		\def\ggga{\alpha}
		\def\gggb{\beta}
		\def\gggc{\gamma}
		\def\gggd{\delta}
		\def\gggp{\phi}
	%........................................................
	% others
	%........................................................
%		\def\half{\frac{1}{2}}
		\def\defeq{\stackrel{\rm def}{=}}


\renewcommand{\baselinestretch}{0.8}


